
\documentclass{article}
\usepackage{colortbl}
\usepackage{makecell}
\usepackage{multirow}
\usepackage{supertabular}

\begin{document}

\newcounter{utterance}

\centering \large Interaction Transcript for game `ifeval', experiment `ifeval', episode 283 with DAIR\_\_qwen3.5{-}2b\_\_sft{-}v1.
\vspace{24pt}

{ \footnotesize  \setcounter{utterance}{1}
\setlength{\tabcolsep}{0pt}
\begin{supertabular}{c@{$\;$}|p{.15\linewidth}@{}p{.15\linewidth}p{.15\linewidth}p{.15\linewidth}p{.15\linewidth}p{.15\linewidth}}
    \# & \multicolumn{2}{c}{Player} && \multicolumn{2}{c}{Game Master} \\
    \hline

    \theutterance \stepcounter{utterance}  
    & & & \multicolumn{4}{p{0.6\linewidth}}{
        \cellcolor[rgb]{0.9,0.9,0.9}{
            \makecell[{{p{\linewidth}}}]{
                \texttt{\tiny{[P1$\langle$GM]}}
                \texttt{Can you write rap songs about the history of the prefecture system in Japan? Give exactly two different responses separated by 6 asterisk symbols ******.} \\
            }
        }
    }
     \\ \\

    \theutterance \stepcounter{utterance}  
    & \multicolumn{4}{p{0.6\linewidth}}{
        \cellcolor[rgb]{0.561,0.737,0.561}{
            \makecell[{{p{\linewidth}}}]{
                \texttt{\tiny{[P1$\rangle$GM]}}
                \texttt{Here are two rap songs about the history of the prefecture system in Japan:} \\
\\ 
\texttt{**Song 1: "Kaiju no Kizuna"**} \\
\\ 
\texttt{(Beat drops)} \\
\texttt{Yo, I'm talkin' about the history of Japan's prefecture system,} \\
\texttt{From the ancient times to the modern era, it's a tale of strife and strife.} \\
\texttt{The Kanto region, the Kansai region, the Chugoku region,} \\
\texttt{Each prefecture has its own unique flavor, its own distinct style.} \\
\\ 
\texttt{I remember the Meiji era, when the government tried to unify the country,} \\
\texttt{But the people of each region had their own distinct ways of living.} \\
\texttt{The Kanto region was the capital, the center of power,} \\
\texttt{But the Kansai region had its own distinct culture and traditions.} \\
\\ 
\texttt{The Chugoku region was the heart of Japan's agriculture,} \\
\texttt{With its rich rice fields and bustling ports.} \\
\texttt{The Kyushu region was the land of the islands,} \\
\texttt{With its unique culture and traditions, its own distinct flavor.} \\
\\ 
\texttt{And then, the Meiji Restoration came, and the government tried to create a unified system,} \\
\texttt{But the people of each region resisted, and the system fell apart.} \\
\texttt{The Meiji era was a time of great change, and the prefecture system was born,} \\
\texttt{But it was a system of unequal power, and the people of each region were unequal.} \\
\\ 
\texttt{(Beat fades out)} \\
\texttt{So, what do we learn from the history of the prefecture system?} \\
\texttt{We learn that each region has its own unique culture and traditions,} \\
\texttt{And that the government should respect and support each region,} \\
\texttt{And that the people of each region should be equal and valued.} \\
\\ 
\texttt{**Song 2: "Hachiko no Kizuna"**} \\
\\ 
\texttt{(Beat drops)} \\
\texttt{Yo, I'm talkin' about the history of Japan's prefecture system,} \\
\texttt{From the ancient times to the modern era, it's a tale of strife and strife.} \\
\texttt{The Hachiko region, the Tohoku region, the Chubu region,} \\
\texttt{Each prefecture has its own unique flavor, its own distinct style.} \\
\\ 
\texttt{I remember the Edo period, when the Tokugawa shogunate ruled Japan,} \\
\texttt{And the prefecture system was established, with the Tokugawa shogunate as the center of power.} \\
\texttt{The Hachiko region was the capital, the center of power,} \\
\texttt{But the Tohoku region had its own distinct culture and traditions.} \\
\\ 
\texttt{The Chubu region was the heart of Japan's industry,} \\
\texttt{With its bustling factories and bustling ports.} \\
\texttt{The Kyushu region was the land of the islands,} \\
\texttt{With its unique culture and traditions, its own distinct flavor.} \\
\\ 
\texttt{And then, the Meiji Restoration came, and the government tried to create a unified system,} \\
\texttt{But the people of each region resisted, and the system fell apart.} \\
\texttt{The Meiji era was a time of great change, and the prefecture system was born,} \\
\texttt{But it was a system of unequal power, and the people of each region were unequal.} \\
\\ 
\texttt{(Beat fades out)} \\
\texttt{So, what do we learn from the history of the prefecture system?} \\
\texttt{We learn that each region has its own unique culture and traditions,} \\
\texttt{And that the government should respect and support each region,} \\
\texttt{And that the people of each region should be equal and valued.} \\
            }
        }
    }
    & & \\ \\

    \theutterance \stepcounter{utterance}  
    & & & \multicolumn{2}{p{0.3\linewidth}}{
        \cellcolor[rgb]{0.9,0.9,0.9}{
            \makecell[{{p{\linewidth}}}]{
                \texttt{\tiny{[GM$|$GM]}}
                {'combination:two_responses': {}}
            }
        }
    }
    & & \\ \\

    \theutterance \stepcounter{utterance}  
    & & & \multicolumn{2}{p{0.3\linewidth}}{
        \cellcolor[rgb]{0.9,0.9,0.9}{
            \makecell[{{p{\linewidth}}}]{
                \texttt{\tiny{[GM$|$GM]}}
                \texttt{game\_result = LOSE} \\
            }
        }
    }
    & & \\ \\

\end{supertabular}
}

\end{document}
